\documentclass[11pt]{article}

\usepackage[margin=0.85in]{geometry}
\usepackage{amsmath,amssymb}
\usepackage{booktabs}
\usepackage{longtable}
\usepackage{array}
\usepackage{enumitem}
\usepackage{xcolor}
\usepackage{hyperref}

\hypersetup{
  colorlinks=true,
  linkcolor=blue!55!black,
  citecolor=blue!55!black,
  urlcolor=blue!55!black
}
\setlist{nosep}
\newcommand{\Hmat}{\mathbf{H}}
\newcommand{\Eideal}{E_{\mathrm{ideal}}}
\newcommand{\code}[1]{\texttt{#1}}

\title{Q-SENSE Error-Mitigation Explorations:\\Methods, Results, and Reproducible Handoff}
\author{}
\date{12 September 2026}

\begin{document}
\maketitle

\begin{abstract}
This document records the error-mitigation studies performed for the Q-SENSE
subspace estimator in this repository.  The main test case was H$_2$O at
$R=1.0$~\AA{} with the optimized variational-orbital (VO) basis and the
\code{FakeQuebec} noise model.  We isolated the dominant matrix-entry error,
verified noiseless finite-shot statistics, tested exact-symmetry control
variates, introduced an exact fixed-orbital sector split, implemented local
readout mitigation, and studied deterministic and probabilistic local-folding
zero-noise extrapolation (ZNE).  We also applied readout mitigation and linear
ZNE to the complete 37-state perturbative (PT) matrix.  The principal findings
are: (i) the original VO energy error is dominated by $H_{00}$; (ii) splitting
the nearly fixed $Z_0$ sector reduces that dominant entry bias by about a factor
of six without approximating the state span; (iii) local readout mitigation is
helpful but does not remove the dominant simulator bias; (iv) curved ZNE fits
reduce bias for one deep split-basis diagonal entry, at substantial variance
cost; and (v) applying the original Q-SENSE shot allocation naively to every
entry and ZNE factor is unreliable for the larger PT matrix.  Numerical claims
below are tied to saved artifacts so they can be checked independently.
\end{abstract}

\section{Scope and status}

All noisy results in this note are local finite-shot Qiskit Aer simulations.
They use the calibration data, connectivity, native gate target, and noise
model supplied by \code{FakeQuebec}; they are not executions on an IBM quantum
processor.  The problem circuits use Aer's \code{statevector} method with
\code{enable\_truncation=True}.  Consequently these are noisy stochastic,
shot-based statevector simulations rather than density-matrix calculations.
For the wide but unentangled readout-calibration circuits used by the
full-matrix scripts, the method was changed to
\code{matrix\_product\_state} to avoid allocating a 127-qubit statevector.

The baseline Q-SENSE accuracy request is
\[
  \epsilon_E=1.6\times10^{-3}\ \mathrm{Ha}.
\]
This value is set by \code{TARGET\_ESTIMATE\_ACCURACY = 1.6e-3} in
\path{scripts/sampling_and_sv.py}.  The original and split VO simulations
predate the portable-output default and therefore retain several trusted local
pickle files.  Newer mitigation runs save numerical data in JSON, NPZ, CSV,
and circuits in QPY wherever practical.

\subsection{Terminology and error metrics}

For a matrix entry, bias means the observed estimate minus its exact noiseless
value,
\[
  b_{ij}=\widehat H_{ij}-H_{ij}^{\mathrm{ideal}}.
\]
For an energy estimate, the same sign convention is used.  Standard
deviations reported by the estimator describe finite-shot uncertainty under
the sampled distribution; they do not include uncertainty in the noise model
or, unless explicitly stated, the readout-calibration matrices.  We use
\[
  \mathrm{RMSE}=\sqrt{b^2+s^2}
\]
when $s^2$ is the estimated sampling variance.  Whenever a noisy matrix
develops spurious eigenvalues below the desired state, both the algebraically
lowest eigenpair and the eigenpair with maximum overlap with the ideal
coefficient vector are reported.  The latter is called the
\emph{overlap-tracked target}; it is a simulator diagnostic because the ideal
vector would not be known on hardware.

\section{Baseline Q-SENSE estimation}

Let $\{\lvert\Phi_i\rangle\}$ be an orthonormal Q-SENSE basis and
\[
  H_{ij}=\langle\Phi_i\rvert\hat H\lvert\Phi_j\rangle,
  \qquad
  \Hmat\vec c=E\vec c.
\]
Each quantum matrix element is decomposed into jointly measurable fragments
$f$.  If $\sigma_{ij,f}$ is the exact single-shot standard deviation of a
fragment, the ideal first-order energy sensitivity gives the weights
\[
  w_{ij}=
  \begin{cases}
    |c_i|^2\sum_f\sigma_{ii,f},&i=j,\\
    2|c_i c_j|\sum_f\sigma_{ij,f},&i>j.
  \end{cases}
\]
Writing $M=\sum_{i\geq j}w_{ij}$, the implemented total-shot estimate and
fragment allocation are, up to ceiling,
\begin{align}
  N&=\left\lceil\frac{M^2}{\epsilon_E^2}\right\rceil,\\
  N_{ij,f}&=\left\lceil
    N\frac{w_{ij}}{M}
    \frac{\sigma_{ij,f}}{\sum_g\sigma_{ij,g}}
  \right\rceil.
\end{align}
Thus observable coefficients matter through the exact fragment standard
deviations and the shot allocation.  They are not encoded in a PUB's
\code{target\_precision}: the local \code{BackendEstimatorV2} chooses its
number of shots as $\lceil1/\epsilon_{\rm PUB}^2\rceil$, so the code passes
\[
  \epsilon_{\rm PUB}=\frac{1+10^{-12}}{\sqrt{N_{ij,f}}}
\]
only to reproduce the already calculated shot count exactly.

Before execution, state-preparation and interference circuits are transpiled
against the \code{FakeQuebec} target with optimization level 3 and transpiler
seed 0.  This enforces its native gate set and limited connectivity.  The
nominal circuit width remains 127 in saved QPY objects, but Aer truncation
simulates only active connected qubits; the completed VO runs reported ranges
of 5--19 active qubits across groups.  A local
\code{CapturingBackendEstimatorV2} is used directly, without a Runtime batch,
and stores the exact measurement histograms used for its expectations.

\subsection{Original 11-state VO baseline}

The completed run is
\begin{verbatim}
saved/runtime_runs/
  20260722T164209Z_h2o_1.0_qwc_fake_quebec/
\end{verbatim}
It used 45 quantum matrix elements, 585 PUBs, and 205,731 retained shots.
The maximum transpiled circuit depth was 1006 and the maximum two-qubit gate
count was 376.  Runtime-level mitigation was disabled.  The result was
\[
  \Eideal=-75.0176664101\ \mathrm{Ha},\qquad
  \widehat E=-74.2180687901\ \mathrm{Ha},
\]
with
\[
 \text{bias}=+0.7995976200\ \mathrm{Ha},\qquad
 \text{standard deviation}=0.7643728997\ \mathrm{Ha},\qquad
 \text{RMSE}=1.1061746172\ \mathrm{Ha}.
\]

The leading ideal coefficient weight is
$|c_0|^2=0.9846882$.  The corresponding diagonal entry had
\[
 H_{00}^{\mathrm{ideal}}=-74.9897246858,\qquad
 \widehat H_{00}=-59.9089790839,
\]
or bias $+15.0807456019$~Ha.  Replacing only this noisy entry by its ideal
value while retaining all other noisy entries gives an energy bias of
$+0.0275449236$~Ha and coefficient-vector fidelity $0.9839845$.  This
counterfactual establishes that $H_{00}$, rather than the accumulation of all
other errors, dominated this particular run.

\subsection{Noiseless check of \texorpdfstring{$H_{00}$}{H00} statistics}

A finite-shot run with no noise model gave entry bias $+0.00246196$~Ha.
A separate 10,000-replicate bootstrap found
\[
\begin{array}{lc}
\toprule
\text{quantity} & \text{standard error (Ha)}\\
\midrule
\text{analytic Q-SENSE prediction} & 0.00139248\\
\text{parametric bootstrap from exact probabilities} & 0.00138930\\
\text{nonparametric bootstrap of the finite histogram} & 0.00156161\\
\bottomrule
\end{array}
\]
The parametric result agrees with the analytic calculation to $0.23\%$.
The larger nonparametric value reflects the particular finite histogram and
rare outcomes.  The saved quantity $0.21810137$ for fragment 0 is a
\emph{single-shot} standard deviation, not the standard error after 60,882
shots.  These checks support the noiseless variance implementation and show
that the large noisy bias is not a failure of ideal shot-noise accounting.

\section{Exact observable shifts as control variates}

Suppose the target state obeys an exact constraint
$\langle C\rangle=c$.  The shifted estimator
\[
  \widehat O_\lambda
  =\widehat O+\lambda(\widehat C-c)
\]
has the same ideal expectation as $O$.  When $O$ and $C$ are evaluated from
the same bit strings,
\[
 \operatorname{Var}(\widehat O_\lambda)
 =\operatorname{Var}(\widehat O)
  +2\lambda\operatorname{Cov}(\widehat O,\widehat C)
  +\lambda^2\operatorname{Var}(\widehat C),
\]
so the variance-optimal coefficient is
\[
 \lambda_{\rm var}=-
 \frac{\operatorname{Cov}(\widehat O,\widehat C)}
      {\operatorname{Var}(\widehat C)}.
\]
This coefficient minimizes variance, not hardware bias.  A coefficient
\[
 \lambda_{\rm bias}=-
 \frac{\widehat O-O_{\rm ideal}}{\widehat C-c}
\]
was also computed to diagnose whether a control could cancel the observed
bias.  It is labelled \emph{diagnostic bias-optimal} because it uses the ideal
answer and must not be selected this way in an unknown production problem.

All results in this section use entry $(0,0)$, fragment 0, its 60,882 saved
shots, and 20,000 multinomial bootstrap resamples.  The ideal fragment value
is $-30.6034385637$~Ha.  The raw bootstrap fragment bias is approximately
$+14.70$~Ha; the other fragments and exact classical contribution make the
raw full-entry bias approximately $+14.80$~Ha.

\subsection{Global parity}

The first control was the tapered-register parity
$P=\prod_q Z_q$ with eigenvalue $p=-1$.  With
$\lambda=-24.8194965$, chosen from the largest observable coefficient, the
sampled parity was $-0.1790020$.  The full-entry bias changed from
$+14.8021$ to $-5.5732$~Ha, while the bootstrap standard deviation increased
from $0.08538$ to $0.12087$~Ha.  The correction therefore crossed the ideal
value: the large parity violation multiplied by a coefficient chosen for
Pauli-coefficient cancellation rather than noise-response calibration.
The variance-optimal coefficient, $-3.2584$, barely improved the bias
($+12.0240$~Ha at fragment level).  The diagnostic bias-optimal coefficient
was $-17.9045$.

\subsection{Pair number}

For seven tapered occupation qubits, the pair-number operator is
\[
  N_{\rm pair}=\frac{7}{2}I-\frac12\sum_{q=0}^{6}Z_q,
  \qquad \langle N_{\rm pair}\rangle=5.
\]
The same noisy bit strings gave $\widehat N_{\rm pair}=4.5888768$, much
closer to the exact value than the parity estimator was to $-1$.  Let
$a=24.8194965$ denote the magnitude used in the coefficient arguments.
The tested results were
\[
\begin{array}{lrrrr}
\toprule
\text{choice} & \lambda & \text{full-entry bias} & \text{bootstrap std.}
              & \text{RMSE}\\
\midrule
\text{raw} & 0 & 14.8014 & 0.08561 & 14.8017\\
2a/5 & 9.9278 & 10.7202 & 0.07320 & 10.7204\\
2a/3 & 16.5463 & 7.9993 & 0.07723 & 7.9997\\
4a/3 & 33.0927 & 1.1972 & 0.11918 & 1.2031\\
\lambda_{\rm var} & 10.6252 & 10.4335 & 0.07314 & 10.4337\\
\bottomrule
\end{array}
\]
The variance-optimal full-entry values were reconstructed from the saved
fragment samples by adding the same fixed contribution as in the other rows.
The diagnostic bias-optimal value was $35.7547$.  The closeness of $2a/5$ to
$\lambda_{\rm var}$ is empirical, not a general identity: $2a/5$ follows a
coefficient-balancing argument, whereas $\lambda_{\rm var}$ follows a sampled
covariance ratio.

\subsection{\texorpdfstring{$Z_0$}{Z0}-weighted number control}

To make the control share more of the dominant $Z_0$ term's noise response,
we tested
\[
  C=(N_{\rm pair}-5)Z_0,
  \qquad \langle C\rangle=0.
\]
For seven qubits,
\[
 C=-\frac12\left(3Z_0+I+\sum_{j\ne0}Z_0Z_j\right).
\]
This algebra explains why $\lambda=2a/3$ cancels the explicit dominant
$Z_0$ coefficient in the shifted observable.  Cancellation of that one
coefficient does \emph{not} force the noisy expectation of the entire shifted
quantum part to vanish, because the identity and six $Z_0Z_j$ terms remain.
The measured control value was $-0.1876745$.

The variance-optimal value $\lambda=10.2948$ reduced the standard deviation
from $0.08559$ to $0.07231$~Ha but left full-entry bias $+12.8710$~Ha.
Coefficient cancellation ($\lambda=16.5463$) left bias $+11.6976$~Ha.
The deliberately large choice $\lambda=2a=49.6390$ reduced full-entry bias to
$+5.4862$~Ha while increasing the standard deviation to $0.18940$~Ha.  The
diagnostic bias-optimal coefficient $78.3250$ produced bias $+0.1019$~Ha and
standard deviation $0.31120$~Ha, but it is target-leaking and therefore not a
usable mitigation prescription.

The control's shared $Z_0$ and $Z_0Z_j$ support gives it a more similar noise
profile to the dominant observable term than $N_{\rm pair}-5$ alone.  This is
a plausible mechanism, not a theorem under the gate-dependent
\code{FakeQuebec} model.  In the actual diagnostic, the unweighted number
control with $\lambda=4a/3$ gave a better operationally motivated bias result
than the tested weighted coefficients.  A future control should be selected
from independent pilot data or a noise model using covariance and response
similarity, then frozen before evaluating production data.

\section{Exact fixed-orbital sector split}

The dominant VO state was found to occupy tapered orbital/qubit 0 almost
entirely in $\lvert1\rangle$.  Rather than use $Z_0$ only as a noisy measured
control, state 0 was split into its normalized fixed-$Z_0$ sectors.  Six pair
generators acting on orbital 0 were removed; the remaining 24 generators were
warm-started from the original angles and independently reoptimized in each
six-qubit active sector.  The results were
\[
\begin{array}{lrr}
\toprule
 & q_0=0 & q_0=1\\
\midrule
\text{sector probability} & 5.1370\times10^{-7} & 0.9999994863\\
\text{number of determinants} & 6 & 15\\
\text{optimized infidelity} & 6.0\times10^{-15} & 1.7\times10^{-13}\\
\text{active-circuit two-qubit gates} & 48 & 48\\
\bottomrule
\end{array}
\]
Their weighted coherent reconstruction has fidelity
$0.9999999999998239$ with the original 21-determinant state.  In the expanded
Q-SENSE basis the two normalized sectors are used as separate states, without
their reconstruction weights.  They are mutually orthogonal by construction,
and the remaining old basis states lie in different pre-existing symmetry
sectors as required by the method.  No generalized eigenvalue problem is
introduced.  Qubit 0 is classical and removed from the active state circuits
for these two states.

The split basis contains 12 rather than 11 states.  Its dominant physical
component is entry $(1,1)$, corresponding to the $q_0=1$ sector, with
$|c_1|^2=0.9846877$.  Its raw entry bias was $+2.5567470$~Ha, versus
$+15.0807456$~Ha for the unsplit $H_{00}$.  If only split entry $(1,1)$ is
set to ideal, the remaining noisy matrix gives energy bias $+0.0228070$~Ha
and coefficient fidelity $0.979334$.  This confirms a large structural bias
reduction for the dominant entry.

The single complete split run nevertheless had energy bias $+0.8839538$~Ha
and coefficient fidelity $1.73\times10^{-5}$ because the biased dominant
diagonal no longer remained the lowest state in the noisy matrix.  It used
51 quantum matrix elements, 626 PUBs, and 201,422 retained shots.  This result
does not negate the entry-level improvement; it shows that the total estimator
still needs mitigation or sufficient sampling for the other entries.

\section{Local readout-error mitigation}

\subsection{Implemented estimator}

For every measured physical qubit $q$, two calibration experiments estimate
\[
 A_q=
 \begin{pmatrix}
 p(0|0)&p(0|1)\\
 p(1|0)&p(1|1)
 \end{pmatrix}.
\]
All support qubits are prepared and measured in parallel, so the independent
local model requires only two calibration circuits: one all-zero and one
all-one circuit, each with $S_{\rm cal}$ shots.  For a measured Pauli product
with support $Q$, the corrected outcome weight is evaluated using
\[
 \bigotimes_{q\in Q} A_q^{-1}.
\]
The code applies these factors directly to each observed bit string and
computes both the corrected first and second moments.  It therefore avoids
constructing a dense $2^{|Q|}\times2^{|Q|}$ matrix.  Calibration cost is
$2S_{\rm cal}$ shots, independent of support width under the parallel local
model, while per-outcome processing is linear in Pauli support.

The statistical overhead is not width-independent.  Defining the readout
contrast
\[
 r_q=1-p(1|0)-p(0|1),
\]
a symmetric-error approximation gives a worst-case variance amplification
for a Pauli product of roughly
\[
 \gamma_Q\simeq\prod_{q\in Q}r_q^{-2}.
\]
Thus local inversion is computationally scalable but can become statistically
expensive for long Pauli strings or poorly conditioned $A_q$.

The preferred split-diagonal calibration used $S_{\rm cal}=20{,}000$ and
physical qubits $\{45,54,62,63,64,65\}$.  The assignment matrices were nearly
symmetric, with single-qubit misassignment probabilities between about
$0.42\%$ and $1.90\%$.  For example,
\[
 A_{62}=\begin{pmatrix}
 0.98105&0.01870\\0.01895&0.98130
 \end{pmatrix}.
\]
This symmetry is why readout twirling was not expected to add a large benefit;
readout twirling itself has not been tested.

\subsection{Split-basis diagonal results}

The paired comparison reran the same circuits and reports raw and mitigated
estimates from the same histograms.  This is the correct comparison; the
``previous bias'' column came from an independent earlier run and can differ
because of finite shots.  For the important $(1,1)$ entry,
\[
  b_{11}^{\rm paired\ raw}=+2.577674\ \mathrm{Ha},\qquad
  b_{11}^{\rm readout}=+2.498453\ \mathrm{Ha}.
\]
Readout inversion removed only $0.07922$~Ha of a roughly $2.58$~Ha bias.
Across the six quantumly sampled diagonals, the absolute paired bias decreased
for all six; the mean absolute bias fell from $1.6487$ to $1.3961$~Ha.
Several of these entries had only 3--317 shots, so individual improvements and
zero empirical standard deviations are noisy diagnostics, not precision
benchmarks.  Fully classical entries are copied exactly and assigned zero
bias and variance by construction.

\subsection{Unpropagated calibration uncertainty}

Current error bars condition on the estimated $A_q$ and omit its finite-shot
uncertainty.  A complete analysis should resample each calibration histogram,
rebuild every $A_q^{-1}$, resample problem histograms, and repeat mitigation,
matrix assembly, and diagonalization.  This hierarchical bootstrap directly
provides the calibration-induced component of entry and energy variance.
Choose $S_{\rm cal}$ so that this component is a small predeclared fraction of
the total variance budget.  The calibration counts are saved, so no new noisy
run is required to implement this analysis.

The reusable code is in \path{src/mitigation.py}; the main switch is
\code{enable\_readout\_mitigation} or environment variable
\code{QSENSE\_ENABLE\_READOUT\_MITIGATION}.  Runtime resilience levels are
not used: the mitigation stack is explicit.

\section{Zero-noise extrapolation}

\subsection{Folding and transpilation convention}

For an eligible native gate $U$, local folding replaces
\[
  U\longmapsto UU^\dagger U,
\]
which is ideal-equivalent but executes additional noisy operations.  Full
local folding was first tested at factors $\{1,3,5\}$.  For a fractional
factor $1\leq\lambda\leq3$, each eligible gate is independently folded with
probability
\[
  p_{\rm fold}=\frac{\lambda-1}{2}.
\]
Eight randomized folding patterns are used for each non-unit factor.  The
factor's fixed shot budget is divided among the patterns; randomization raises
the circuit/PUB count, not the total shots, unless the user separately raises
the per-factor budget.

The source circuits have already been transpiled at optimization level 3 to
the \code{FakeQuebec} native basis and layout.  Folding is then performed in
that basis.  Only operations whose backend target reports nonzero error are
eligible.  IBM-style $R_z$ operations are normally virtual frame changes and
the fake backend reports zero error, so they are not folded.  This also
typically holds on a real IBM device when the final ISA circuit retains an
$R_z$, but the live backend target must be checked.  The inverse of an
\code{sx} is implemented in the native basis as
$R_z(\pi)\,\code{sx}\,R_z(-\pi)$.  If folding introduces an unsupported gate,
it is transpiled at optimization level 0 with the original layout to avoid
optimizing away the fold.

The estimator constructs measurement-basis rotations after this folding
stage.  Those rotations are not folded in the current implementation.
Therefore $\lambda$ scales the saved state/interference circuit's eligible
gates, not every error source in the complete measured circuit.  Realized
gate counts and fold seeds are saved and should be checked.

In an ideal noiseless calculation all folding patterns implement the same
unitary, so their exact expectation values are identical.  Under
gate-dependent and location-dependent noise, patterns with the same expected
$\lambda$ can have different means.  Finite-shot noise can mask this
between-pattern variance; the implementation combines resolved pattern spread
with within-pattern shot variance.

\subsection{Linear and variance-weighted extrapolation}

For factors $\lambda_i$, ordinary linear ZNE fits
\[
  E(\lambda)=\beta_0+\beta_1\lambda,
\]
and takes $\widehat E(0)=\widehat\beta_0=\sum_i c_i\widehat E_i$.
Variance-weighted linear ZNE minimizes
\[
  \chi^2=\sum_i
  \frac{[\widehat E_i-(\beta_0+\beta_1\lambda_i)]^2}{v_i},
\]
using inverse total factor variances.  For independent factors,
\[
 \operatorname{Var}[\widehat E(0)]=\sum_i c_i^2v_i.
\]
If the single-shot standard deviations are $\sigma_i$ and the coefficients
are fixed, the variance-optimal allocation at fixed total shots is
\[
  N_i\propto |c_i|\sigma_i.
\]
There is a design loop: changing $N_i$ changes $v_i$, which changes the
weighted-fit coefficients $c_i$.  The consistent practical solution is to
estimate $\sigma_i$ and weights in an independent pilot, freeze them, impose a
minimum diagnostic allocation at interior factors, and then execute the
production allocation once.  Repeatedly adapting weights to the same data
would invalidate a simple uncertainty interpretation.

\subsection{Full-fold diagnostic on split entry
\texorpdfstring{$(0,0)$}{(0,0)}}

The initial $\{1,3,5\}$ full-fold test concerned the low-weight split entry
$(0,0)$, not the dominant $(1,1)$ entry.  With 60,000 shots per factor and
readout mitigation, the factor estimates were
\[
 (-27.66718,-23.69699,-21.66888)\ \mathrm{Ha}.
\]
Linear extrapolation gave $-28.84307$~Ha versus ideal $-31.03017$~Ha:
bias $+2.18710$~Ha, empirical standard deviation $0.22192$~Ha, and RMSE
$2.19833$~Ha.  The scale-one bias was $+3.36299$~Ha.  This result is distinct
from the earlier $2.5777\rightarrow2.4985$~Ha readout result, which concerned
split entry $(1,1)$ and did not apply ZNE.

Factors 3 and 5 were judged unnecessarily severe for these already deep,
noisy circuits because amplified values moved toward saturation.  Subsequent
work used smaller probabilistic factors.

\subsection{Probabilistic ZNE on dominant split entry
\texorpdfstring{$(1,1)$}{(1,1)}}

With factors $\{1,1.5,2\}$, eight patterns per non-unit factor, and the cached
per-factor fragment allocation, the scale-one and linear-ZNE results were
\[
\begin{array}{lrrr}
\toprule
 & \text{bias (Ha)} & \text{std. (Ha)} & \text{RMSE (Ha)}\\
\midrule
\lambda=1 & 2.49492 & 0.06567 & 2.49579\\
\text{ordinary linear ZNE} & 1.30222 & 0.14826 & 1.31063\\
\bottomrule
\end{array}
\]
The extrapolation reduced bias by about $48\%$ but more than doubled the
standard deviation.

A pilot-based allocation preserved the same total of 441,951 shots but used
\[
 (N_1,N_{1.5},N_2)=(234371,45356,162224).
\]
The fitted intercept coefficients became
\[
 (c_1,c_{1.5},c_2)=(1.98889,0.02221,-1.01111).
\]
The resulting bias, standard deviation, and RMSE were $1.21341$, $0.12433$,
and $1.21976$~Ha, respectively.  Relative to the equal-allocation run, the
variance decreased by $29.7\%$ and RMSE by $6.9\%$.

\subsection{Curved models and five-factor study}

The final grid was
\[
 \lambda\in\{1,1.25,1.5,1.75,2\}.
\]
Two independently seeded runs were inverse-variance pooled, for a nominal
883,902 shots per factor.  The pooled data were
\[
\begin{array}{crr}
\toprule
\lambda&\widehat E(\lambda)&\text{standard deviation}\\
\midrule
1.00&-72.509300&0.02681\\
1.25&-72.133633&0.02773\\
1.50&-71.774107&0.03487\\
1.75&-71.453188&0.02890\\
2.00&-71.160255&0.03185\\
\bottomrule
\end{array}
\]
Three variance-weighted extrapolators were compared:
\begin{align}
 E_{\rm lin}(\lambda)&=\beta_0+\beta_1\lambda,\\
 E_{\rm quad}(\lambda)&=\beta_0+\beta_1\lambda+\beta_2\lambda^2,\\
 E_{\rm exp}(\lambda)&=E_0+s[1-e^{-k\lambda}]/k.
\end{align}
Factor-level parametric bootstrap with 10,000 samples gave
\[
\begin{array}{lrrrr}
\toprule
\text{model}&E(0)&\text{bias}&\text{std.}&\text{RMSE}\\
\midrule
\text{linear}&-73.84512&1.14458&0.0554&1.1459\\
\text{quadratic}&-74.32102&0.66868&0.2752&0.7231\\
\text{offset exponential}&-74.40536&0.58434&0.3809&0.6975\\
\bottomrule
\end{array}
\]
Curved models reduced observed simulator bias but amplified variance.  The
quadratic coefficient was $\beta_2=-0.2294\pm0.1293$ with bootstrap 95\%
interval $[-0.4849,0.0254]$ and two-sided $p=0.076$.  Holding out
$\lambda=1.75$, standardized prediction residuals were 0.632, 0.041, and
0.087 for linear, quadratic, and exponential fits.  Curvature is therefore
supported by the trend and prediction test but was not conventionally
significant at $p<0.05$.

For Gaussian residuals, the reported
$\mathrm{AIC}=\chi^2+2k$ penalizes fit residuals by twice the number of model
parameters.  The AIC values were 7.129, 6.034, and 6.051.  Differences below
roughly 2 do not decisively select a model, so the quadratic's 1.095 advantage
over linear was suggestive rather than conclusive.  More independent factors
are more valuable for detecting curvature than simply reducing error bars on
the same three points, although both help.  There is no universal circuit
depth at which a curved model becomes correct; accumulated error, observable
support, gate composition, and saturation matter together
\cite{majumdar2023,cai2021}.

\section{Complete-matrix mitigation experiments}

\subsection{Original 11-state VO matrix}

A script was written to fix $(0,0)$ at its ideal value and apply local readout
mitigation plus probabilistic linear ZNE to every other quantum entry:
\path{scripts/run_full_matrix_readout_linear_zne.py}.  The first attempt
failed because readout calibration with a 127-qubit statevector exceeded
\code{max\_memory\_mb}.  Calibration was changed to MPS, while problem
circuits retained truncated noisy statevector simulation.  No complete
11-state full-matrix mitigated result is currently saved.  The incomplete
directories ending in \code{20260813T032929Z} and
\code{20260813T033440Z} must not be treated as results.

The useful unmitigated counterfactual is the calculation above: with ideal
$H_{00}$ and every other original noisy entry unchanged, the energy bias is
$+0.0275449$~Ha.  This is the baseline a future completed full-matrix
mitigation run should beat.

\subsection{Completed 37-state PT matrix}

The analogous PT workflow completed and is stored under
\begin{verbatim}
saved/runtime_runs/pt_h2o_1.0/
  pt_readout_linear_zne_20260813T134200Z/
\end{verbatim}
It fixed $(0,0)$ to ideal, used factors $\{1,1.5,2\}$, eight patterns per
non-unit factor, and respective totals of 19,900, 33,937, and 33,937 shots.
There were 40,358 folded PUBs in 17 sequential simulator batches.  The run
used the original Q-SENSE allocation independently at each factor.

\[
\begin{array}{lrrrr}
\toprule
\text{matrix}&\text{lowest bias}&\text{lowest overlap}
 &\text{tracked bias}&\text{tracked overlap}\\
\midrule
\text{same-shot raw}&-8.35454&0.000781&-0.005163&0.91069\\
\text{readout only}&-9.17104&0.000777&-0.068029&0.71424\\
\text{readout + linear ZNE}&-34.28791&0.000437&-0.113953&0.80240\\
\bottomrule
\end{array}
\]
Large noisy diagonal errors create spurious low eigenvalues, so the lowest
energy is not a meaningful estimate here.  Even for the overlap-tracked
state, readout correction and linear ZNE worsened energy bias relative to a
fortuitously accurate raw result.  This is evidence against applying the
baseline allocation mechanically to all entries and factors; it is not
evidence that ZNE must fail after a redesigned allocation.

Many fragments received one or only a few shots.  A one-shot empirical
variance is zero but does not mean the estimator is exact.  Post-processing
therefore used the floor
\[
 v_{ij}^{\rm used}=\max\left(
 v_{ij}^{\rm empirical},
 \sum_f\frac{\sigma_{ij,f,\rm ideal}^2}{N_{ij,f}}
 \right).
\]
This floor was applied 506 times.  Forty-seven ill-conditioned weighted
linear fits fell back to ordinary least-squares coefficients while retaining
the same linear model.  These counts diagnose an underresolved design.

\subsection{Leading PT fragment and increased-shot test}

Entry $(23,0)$, fragment 0 was isolated as a leading cause of coefficient
vector error.  Its ideal entry is $2.3553\times10^{-5}$~Ha, but the original
four-shot raw fragment estimate was $-0.17943$~Ha.  The transpiled fragment
circuit has depth 254 and 90 native ECR gates.  The two PT states contain only
two and one pair rotations, with angles $-6.99\times10^{-4}$ and
$-4.78\times10^{-4}$ radians; their individual logical preparations have
depths 15 and 8 with four and three CX gates.  Hence most of the 90-ECR cost
comes from the interference/controlled-swap construction and routing, not
from the small pair rotations.

Eight independent noisy replicas at 100, 500, 2,000, and 10,000 shots showed
that more shots reduce spread but do not remove device-model bias.  At 10,000
shots the mean entry biases were $+0.03780$~Ha raw and $+0.04215$~Ha after
readout correction.  The corresponding overlap-tracked energy biases were
$+0.037685$ and $+0.032139$~Ha.  Thus the original four-shot value happened
to cancel other matrix errors and gave a misleadingly small total-energy
bias.  In a separate noise-free study with 32 replicas per shot count, the
mean bias remained statistically consistent with zero and the single-run
standard deviation decreased from $3.03\times10^{-4}$ at 100 shots to
$3.36\times10^{-5}$ at 10,000 shots, consistent with $N^{-1/2}$ convergence.

\section{Conclusions supported by the data}

\begin{enumerate}
  \item The original H$_2$O VO error at $R=1.0$~\AA{} is overwhelmingly tied
        to one high-weight diagonal entry.  Increasing shots cannot remove
        its noise-model bias.
  \item Exact control variates can trade bias against variance when the
        control experiences correlated noise.  Variance-optimal and
        bias-optimal coefficients answer different questions.  Coefficients
        fitted using the ideal target are diagnostics only.
  \item The fixed-$Z_0$ sector split is the strongest presently demonstrated
        operational improvement to the dominant entry.  It changes the basis
        and measurement cost but preserves the original state exactly within
        numerical optimization error.
  \item Local readout inversion consistently improved the paired sampled
        diagonals, but the dominant-entry improvement was modest.  Readout is
        therefore not the main error source under this fake-backend model.
  \item Probabilistic ZNE with small factors is preferable here to full
        factors 3 and 5.  For split entry $(1,1)$, ZNE roughly halved bias.
        Curved models reduced RMSE further, but with much larger uncertainty
        and no decisive model-selection result.
  \item The completed PT matrix shows that entrywise mitigation can destabilize
        diagonalization when shots are allocated for the unmitigated
        first-order energy objective.  Full-matrix mitigation needs a new
        joint allocation accounting for extrapolation coefficients, readout
        overhead, minimum per-fragment resolution, and eigenvalue sensitivity.
\end{enumerate}

\section{Known limitations and recommended next work}

\begin{enumerate}
  \item Implement a hierarchical bootstrap that resamples calibration and
        problem histograms together.  Use it to choose $S_{\rm cal}$ and to
        propagate uncertainty through diagonalization.
  \item Redesign full-matrix allocation using pilot estimates of each
        factor's variance and fixed extrapolation coefficients.  Enforce a
        scientifically chosen minimum number of shots per retained fragment;
        never interpret one-shot empirical variance as zero uncertainty.
  \item Fold the complete physical measurement circuit, including estimator
        basis rotations, or explicitly measure the unamplified contribution.
  \item Test model stability under modest changes of factor grid, fold seed,
        and held-out factor.  Reject extrapolation when amplified values have
        saturated or the intercept is unstable.
  \item Explore additional exact controls that share the dominant Pauli
        support.  Select coefficients from independent pilots or a justified
        noise model, not proximity to an ideal answer.
  \item Generalize the fixed-orbital split only where occupation diagnostics
        show an almost fixed tapered qubit.  Record reconstruction fidelity,
        enlarged matrix cost, and resulting shot allocation at every geometry.
  \item Revalidate the conclusions on current real-device calibration data.
        \code{FakeQuebec} establishes behavior of one local model, not a
        hardware-independent result.
\end{enumerate}

\section{Reproduction guide}

Run commands from the repository's inner \code{seniority/} directory using
the \code{SenDev} environment.  The principal scripts are:
\begin{verbatim}
# Build and verify the exact fixed-orbital split
python scripts/build_orbital0_split_states.py

# Baseline split sampling and local noisy simulation
python scripts/sampling_and_sv_orbital0_split.py

# Readout-mitigated split diagonals
python scripts/run_readout_mitigated_diagonals.py \
  --run-dir saved/runtime_runs/<split-run>

# One-entry probabilistic readout + ZNE
python scripts/run_probabilistic_zne_entry.py --entry 1

# Compare linear, quadratic, and exponential models
python scripts/fit_zne_models.py <result-directories>

# Inspect a full-matrix job without running it
python scripts/run_full_matrix_readout_linear_zne.py --dry-run

# Resume PT aggregation from already completed batches
python scripts/run_pt_readout_linear_zne.py \
  --resume-run saved/runtime_runs/pt_h2o_1.0/\
pt_readout_linear_zne_20260813T134200Z
\end{verbatim}
Command-line help should be consulted before a new expensive run because
defaults and required paths differ among the exploratory scripts.

Portable JSON can be loaded without reconstructing Qiskit result objects:
\begin{verbatim}
from pathlib import Path
from src.saved_run_io import load_json
import numpy as np

run = Path("saved/runtime_runs/<run-directory>")
metadata = load_json(run / "run_results.json.gz")
histograms = load_json(run / "raw_measurement_histograms.json.gz")
arrays = np.load(run / "matrices_and_eigenvectors.npz",
                 allow_pickle=False)

ideal_matrix = arrays["ideal_matrix"]
noisy_matrix = arrays["noisy_matrix"]
ideal_c = arrays["ideal_eigenvector"]
noisy_c = arrays["noisy_eigenvector"]
\end{verbatim}
For mitigation result directories, the principal files are
\begin{verbatim}
results.json.gz
fragment_results.json.gz
batches/*.json.gz
\end{verbatim}
as indicated by the script README.  Load QPY circuits with
\code{qiskit.qpy.load}.  Legacy \code{*.pkl} and \code{*.dump} files should
be unpickled only when they are trusted local artifacts; Qiskit container
pickles are version-sensitive and one historical raw \code{DataBin} pickle is
known to be unreadable under a newer Qiskit version.

\subsection{Authoritative artifact map}

\begin{longtable}{>{\raggedright\arraybackslash}p{0.31\textwidth}
                  >{\raggedright\arraybackslash}p{0.61\textwidth}}
\toprule
\text{Purpose} & \text{Repository-relative path}\\
\midrule
\endhead
Original VO baseline &
\path{saved/runtime_runs/20260722T164209Z_h2o_1.0_qwc_fake_quebec/}\\
Noiseless and symmetry-shift diagnostics &
Same directory; files beginning \path{bootstrap_noiseless},
\path{noiseless_control}, \path{symmetry_shift}, \path{number_shift}, and
\path{z0_weighted_number_shift}.\\
Split-state construction &
\path{data/h2o_data/UCSF_sym_comp_for_Praveen_Smik_1.0_state0_orbital0_split.dump}
and companion \path{.txt}.\\
Split VO baseline &
\path{saved/runtime_runs/20260724T024719Z_h2o_1.0_state0_orbital0_split_v1_qwc_fake_quebec/}\\
Preferred readout diagonal run &
Split run plus \path{readout_mitigated_diagonals_20260724T211017Z/}.\\
Initial and optimized three-factor ZNE &
Split run plus \path{probabilistic_zne_entry_1_1_20260725T200028Z/} and
\path{...20260725T203218Z/}.\\
Final five-factor studies &
Split run plus \path{probabilistic_zne_entry_1_1_20260727T161644Z/} and
\path{...20260727T162129Z/}.\\
Detailed ZNE-only note & \path{documents/zne_study_summary.tex}.\\
Completed PT full-matrix run &
\path{saved/runtime_runs/pt_h2o_1.0/pt_readout_linear_zne_20260813T134200Z/}.\\
Reusable mitigation routines & \path{src/mitigation.py}.\\
Raw-count estimator adapter & \path{src/capturing_backend_estimator.py}.\\
Script inventory & \path{scripts/README.md}.\\
\bottomrule
\end{longtable}

\begin{thebibliography}{9}

\bibitem{giurgica2020}
T.~Giurgica-Tiron, Y.~Hindy, R.~LaRose, A.~Mari, and W.~J.~Zeng,
``Digital zero noise extrapolation for quantum error mitigation,''
\emph{2020 IEEE International Conference on Quantum Computing and
Engineering}, 306--316 (2020).
\href{https://doi.org/10.1109/QCE49297.2020.00045}
{doi:10.1109/QCE49297.2020.00045}.

\bibitem{majumdar2023}
R.~Majumdar, P.~Rivero, F.~Metz, A.~Hasan, and D.~S.~Wang,
``Best practices for quantum error mitigation with digital zero-noise
extrapolation,'' \emph{2023 IEEE International Conference on Quantum
Computing and Engineering}, 881--887 (2023).
\href{https://doi.org/10.1109/QCE57702.2023.00102}
{doi:10.1109/QCE57702.2023.00102}.

\bibitem{cai2021}
Z.~Cai,
``Multi-exponential error extrapolation and combining error mitigation
techniques for NISQ applications,'' \emph{npj Quantum Information}
\textbf{7}, 80 (2021).
\href{https://doi.org/10.1038/s41534-021-00404-3}
{doi:10.1038/s41534-021-00404-3}.

\bibitem{krebsbach2022}
M.~Krebsbach, B.~Trauzettel, and A.~Calzona,
``Optimization of Richardson extrapolation for quantum error mitigation,''
arXiv:2201.08080 (2022). \url{https://arxiv.org/abs/2201.08080}.

\end{thebibliography}

\end{document}
